\documentclass[10pt,a4paper]{amsart}
\usepackage{amsmath,amssymb,amsthm,graphicx}
% Retained minimal preamble for the Stacks Project chapter flat.tex (More on Flatness).
% Verbatim excerpt of the chapter preamble: Xy-pic setup (native preamble lines 14-18) and the native
% theorem declarations with the mathematical macros (native preamble lines 154-217). Copied mechanically;
% nothing added, removed or redefined. The wrapper class and the amsmath/amssymb/amsthm/graphicx base
% packages are supplied by the assembler, so the class/externaldocument part of the native preamble is omitted.
% For commutative diagrams we use Xy-pic
\usepackage[all]{xy}

% We use 2cell for 2-commutative diagrams.
\xyoption{2cell}
% Theorem environments.
%
\theoremstyle{plain}
\newtheorem{theorem}[subsection]{Theorem}
\newtheorem{proposition}[subsection]{Proposition}
\newtheorem{lemma}[subsection]{Lemma}

\theoremstyle{definition}
\newtheorem{definition}[subsection]{Definition}
\newtheorem{example}[subsection]{Example}
\newtheorem{exercise}[subsection]{Exercise}
\newtheorem{situation}[subsection]{Situation}

\theoremstyle{remark}
\newtheorem{remark}[subsection]{Remark}
\newtheorem{remarks}[subsection]{Remarks}

\numberwithin{equation}{subsection}

% Macros
%
\def\lim{\mathop{\mathrm{lim}}\nolimits}
\def\colim{\mathop{\mathrm{colim}}\nolimits}
\def\Spec{\mathop{\mathrm{Spec}}}
\def\Hom{\mathop{\mathrm{Hom}}\nolimits}
\def\Ext{\mathop{\mathrm{Ext}}\nolimits}
\def\SheafHom{\mathop{\mathcal{H}\!\mathit{om}}\nolimits}
\def\SheafExt{\mathop{\mathcal{E}\!\mathit{xt}}\nolimits}
\def\Sch{\mathit{Sch}}
\def\Mor{\mathop{\mathrm{Mor}}\nolimits}
\def\Ob{\mathop{\mathrm{Ob}}\nolimits}
\def\Sh{\mathop{\mathit{Sh}}\nolimits}
\def\NL{\mathop{N\!L}\nolimits}
\def\CH{\mathop{\mathrm{CH}}\nolimits}
\def\proetale{{pro\text{-}\acute{e}tale}}
\def\etale{{\acute{e}tale}}
\def\QCoh{\mathit{QCoh}}
\def\Ker{\mathop{\mathrm{Ker}}}
\def\Im{\mathop{\mathrm{Im}}}
\def\Coker{\mathop{\mathrm{Coker}}}
\def\Coim{\mathop{\mathrm{Coim}}}

% Boxtimes
%
\DeclareMathSymbol{\boxtimes}{\mathbin}{AMSa}{"02}

%
% Macros for moduli stacks/spaces
%
\def\QCohstack{\mathcal{QC}\!\mathit{oh}}
\def\Cohstack{\mathcal{C}\!\mathit{oh}}
\def\Spacesstack{\mathcal{S}\!\mathit{paces}}
\def\Quotfunctor{\mathrm{Quot}}
\def\Hilbfunctor{\mathrm{Hilb}}
\def\Curvesstack{\mathcal{C}\!\mathit{urves}}
\def\Polarizedstack{\mathcal{P}\!\mathit{olarized}}
\def\Complexesstack{\mathcal{C}\!\mathit{omplexes}}
% \Pic is the operator that assigns to X its picard group, usage \Pic(X)
% \Picardstack_{X/B} denotes the Picard stack of X over B
% \Picardfunctor_{X/B} denotes the Picard functor of X over B
\def\Pic{\mathop{\mathrm{Pic}}\nolimits}
\def\Picardstack{\mathcal{P}\!\mathit{ic}}
\def\Picardfunctor{\mathrm{Pic}}
\def\Deformationcategory{\mathcal{D}\!\mathit{ef}}
\usepackage{cleveref}
\crefname{originallemma}{original lemma}{original lemmas}
\Crefname{originallemma}{Original lemma}{Original lemmas}
\makeatletter
\newlabel{divisors-lemma-blowing-up-two-ideals}{{080A}{source}{}{sourcepointer-1}{}}
\newlabel{divisors-lemma-blowing-up-two-ideals@cref}{{[originallemma][0][]080A}{[1][0][]source}{}{}{}}
\makeatother
\crefname{originallemma}{original lemma}{original lemmas}
\Crefname{originallemma}{Original lemma}{Original lemmas}
\makeatletter
\newlabel{algebra-lemma-add-trivial-complex}{{00MT}{source}{}{sourcepointer-2}{}}
\newlabel{algebra-lemma-add-trivial-complex@cref}{{[originallemma][0][]00MT}{[1][0][]source}{}{}{}}
\makeatother
\title{Existence of a U-admissible blowup on which a perfect module of tor dimension at most one vanishing on U acquires two canonical filtrations with prescribed invertible Fitting divisors}
\author{The Stacks Project Authors}
\date{stacks-project:flat.tex:local-snapshot-20261011; GNU FDL 1.2 or later}
\begin{document}
\maketitle
\noindent\textit{Editorial note.} All delivered mathematical text is the author's own, copied from the pinned Stacks Project chapter file flat.tex; no statement or proof step was rewritten, reordered or removed.\par\medskip
\noindent\textit{Editorial note.} The preamble is the chapter's own retained minimal preamble (sources/stacks-flat/retained-preamble.tex), a verbatim excerpt of the native Xy-pic setup and of the native theorem declarations and mathematical macros; no macro was invented or redefined, and the wrapper class and the amsmath/amssymb/amsthm/graphicx base packages are supplied by the assembler.\par\medskip
\noindent\textit{Editorial note.} The counted statement and proof contain no cross reference, no citation and no included file, so the delivered body is self contained.\par\medskip
\noindent\textit{Editorial note.} Canonical location of the counted proposition: Stacks tag 0ESU (sources/stacks-flat/pinned-original-source.tex line 13266).\par\medskip
\noindent\textit{Allowed standard prerequisites (editorial).}\begin{itemize}
\item Fitting ideals of a finitely presented module and their behaviour under blowing up two ideals. Divisors, Lemma labelled divisors-lemma-blowing-up-two-ideals (Stacks tag 080A), invoked at proof lines 13301-13302. Supplies that after blowing up the product of the Fitting ideals they become invertible; the definition of the blowup used and the reduction to the invertible case are carried out in the counted span.
\item Dropping a trivial summand from a presentation of a module. Algebra, Lemma labelled algebra-lemma-add-trivial-complex (Stacks tag 00MT), invoked at proof line 13338. Supplies that a presentation whose matrix has entries generating the unit ideal can be replaced by one with n decreased by one, which is the induction step on n; the construction of the divisors and of the two filtrations stays in the counted span.
\end{itemize}\medskip
\ifdefined\hypertarget\hypertarget{sourcepointer-1}{}\fi
\noindent\textit{Source reference.} Original lemma 080A: sources/stacks-chow/tags.txt, source lines 10175--10175. Cross-chapter standard prerequisite located in the project tag list staged in the local snapshot; invoked at proof line 13301 of the counted proof. It is retained only as a navigation pointer; every load-bearing step is inside the delivered span.\par\medskip
\ifdefined\hypertarget\hypertarget{sourcepointer-2}{}\fi
\noindent\textit{Source reference.} Original lemma 00MT: sources/stacks-chow/tags.txt, source lines 835--835. Cross-chapter standard prerequisite located in the project tag list staged in the local snapshot; invoked at proof line 13338 of the counted proof. It is retained only as a navigation pointer; every load-bearing step is inside the delivered span.\par\medskip
\begin{lemma}
\label{lemma-canonical-blowup-torsion-pd1}
Let $X$ be a scheme. Let $\mathcal{F}$ be a perfect $\mathcal{O}_X$-module
of tor dimension $\leq 1$. Let $U \subset X$ be an open such that
$\mathcal{F}|_U = 0$. Then there is a $U$-admissible blowup
$$
b : X' \to X
$$
such that $\mathcal{F}' = b^*\mathcal{F}$ is equipped with two canonical
locally finite filtrations
$$
0 = F^0 \subset F^1 \subset F^2 \subset \ldots \subset \mathcal{F}'
\quad\text{and}\quad
\mathcal{F}' = F_1 \supset F_2 \supset F_3 \supset \ldots \supset 0
$$
such that for each $n \geq 1$ there is an effective Cartier divisor
$D_n \subset X'$ with the property that
$$
F^i/F^{i - 1}
\quad\text{and}\quad
F_i/F_{i + 1}
$$
are finite locally free of rank $i$ on $D_i$.
\end{lemma}

\begin{proof}
Choose an affine open $V \subset X$ such that there exists a presentation
$$
0 \to \mathcal{O}_V^{\oplus n} \xrightarrow{A} \mathcal{O}_V^{\oplus n} \to
\mathcal{F} \to 0
$$
for some $n$ and some matrix $A$. The ideal we are going to blowup in
is the product of the Fitting ideals $\text{Fit}_k(\mathcal{F})$
for $k \geq 0$. This makes sense because in the affine
situation above we see that $\text{Fit}_k(\mathcal{F})|_V = \mathcal{O}_V$
for $k > n$. It is clear that this is a $U$-admissible blowing up. By
Divisors, Lemma \ref{divisors-lemma-blowing-up-two-ideals}
we see that on $X'$ the ideals $\text{Fit}_k(\mathcal{F})$
are invertible. Thus we reduce to the case discussed in the
next paragraph.

\medskip\noindent
Assume $\text{Fit}_k(\mathcal{F})$ is an invertible ideal for
$k \geq 0$. If $E_k \subset X$ is the effective
Cartier divisor defined by $\text{Fit}_k(\mathcal{F})$
for $k \geq 0$, then the effective Cartier divisors
$D_k$ in the statement of the lemma will satisfy
$$
E_k = D_{k + 1} + 2 D_{k + 2} + 3 D_{k + 3} + \ldots
$$
This makes sense as the collection $D_k$ will be locally finite.
Moreover, it uniquely determines the effective Cartier divisors $D_k$
hence it suffices to construct $D_k$ locally.

\medskip\noindent
Choose an affine open $V \subset X$ and presentation of $\mathcal{F}|_V$
as above. We will construct the divisors and filtrations by induction
on the integer $n$ in the presentation. We set $D_k|_V = \emptyset$
for $k > n$ and we set $D_n|V = E_{n - 1}|_V$.
After shrinking $V$ we may assume that
$\text{Fit}_{n - 1}(\mathcal{F})|_V$ is generated by a
single nonzerodivisor $f \in \Gamma(V, \mathcal{O}_V)$.
Since $\text{Fit}_{n - 1}(\mathcal{F})|_V$ is the ideal generated
by the entries of $A$, we see that there is a matrix
$A'$ in $\Gamma(V, \mathcal{O}_V)$ such that $A = fA'$.
Define $\mathcal{F}'$ on $V$ by the short exact sequence
$$
0 \to \mathcal{O}_V^{\oplus n} \xrightarrow{A'} \mathcal{O}_V^{\oplus n} \to
\mathcal{F}' \to 0
$$
Since the entries of $A'$ generate the unit ideal in
$\Gamma(V, \mathcal{O}_V)$ we see that $\mathcal{F}'$
locally on $V$ has a presentation with $n$ decreased
by $1$, see Algebra, Lemma \ref{algebra-lemma-add-trivial-complex}.
Further note that
$f^{n - k}\text{Fit}_k(\mathcal{F}') = \text{Fit}_k(\mathcal{F})|_V$
for $k = 0, \ldots, n$. Hence $\text{Fit}_k(\mathcal{F}')$
is an invertible ideal for all $k$. We conclude by induction
that there exist effective Cartier divisors $D'_k \subset V$
such that $\mathcal{F}'$ has two canonical filtrations as in the statement
of the lemma. Then we set $D_k|_V = D'_k$ for $k = 1, \ldots, n - 1$.
Observe that the equalities between effective Cartier divisors
displayed above hold with these choices. Finally, we come
to the construction of the filtrations. Namely,
we have short exact sequences
$$
0 \to
\mathcal{O}_{D_n \cap V}^{\oplus n} \to
\mathcal{F} \to \mathcal{F}' \to 0
\quad\text{and}\quad
0 \to
\mathcal{F}' \to \mathcal{F} \to
\mathcal{O}_{D_n \cap V}^{\oplus n} \to 0
$$
coming from the two factorizations $A = A'f = f A'$ of $A$.
These sequences are canonical because in the first one
the submodule is $\Ker(f : \mathcal{F} \to \mathcal{F})$
and in the second one the quotient module is
$\Coker(f : \mathcal{F} \to \mathcal{F})$.
\end{proof}

\end{document}
